\documentclass[10pt,a4paper]{article}
\usepackage[T1]{fontenc}
\usepackage[margin=0.58in]{geometry}
\usepackage{enumitem}
\usepackage{booktabs}
\setlist[itemize]{leftmargin=1.0em,nosep,topsep=1pt,parsep=0pt}
\renewcommand{\arraystretch}{1.02}
\pagestyle{empty}

\begin{document}
\footnotesize

\begin{center}
{\large\bfseries Decision Log}\\[2pt]
H1 \texttt{provider\_services\_agreement.md} (INS-H1-2024-0417) \quad
H4 \texttt{conditional\_reimbursement\_agreement.md} (INS-H4-2024-2049) \quad
H5 \texttt{network\_reimbursement\_agreement.md} (INS-H5-2024-0731)\\[2pt]
Hospital 1 is the labelled development set (calibration only, not scored); Hospitals 4 and 5 are scored.
Hospitals 2 and 3 were not attempted.
\end{center}

\vspace{-3pt}
\section*{\normalsize 1.\; Assumptions}
\vspace{-4pt}

\begin{center}
\begin{tabular}{@{}llll@{}}
\toprule
& \textbf{Hospital 1} & \textbf{Hospital 4} & \textbf{Hospital 5} \\
\midrule
Base rates            & Section 4              & Section 3               & Section 4, Table 1 \\
Daily caps            & Section 8              & Section 6               & column of Table 1 \\
Facility multiplier   & single site, $1\times$ & single site, $1\times$  & Table 2, 3 sites \\
Plan-tier multiplier  & tier-blind, $1\times$  & tier-blind, $1\times$   & Table 3, 3 tiers \\
Weekend uplifts       & Section 6              & Section 10: ``None.''   & Section 6 \\
Invoicing checks      & Section 11             & Section 11              & Section 10 \\
\bottomrule
\end{tabular}
\end{center}
\vspace{-2pt}

\begin{itemize}
\item \textbf{Execution order} (H1 3.2, H4 4.1, H5 3.1), identical in all three contracts: bundle
substitution $\to$ facility multiplier $\to$ plan-tier multiplier $\to$ premium/uplift $\to$ cumulative
volume discount; the line total is the resulting unit rate $\times$ billed quantity.
\item \textbf{Rounding:} integer cents, half-up, applied after \emph{each} step rather than once at the end.
Material, not cosmetic: through Hospital 5's two multiplier stages alone the conventions diverge on roughly
7\% of sampled lines (e.g.\ $156{,}675\times1.1\times0.92$ gives 158{,}556 step-wise vs 158{,}555 rounded once).
\item \textbf{Structure is re-derived per contract}, never aligned positionally --- section numbering differs
across all three. Every table is parsed from its own clause heading and its row count asserted against a
manual read before anything downstream uses it.
\item \textbf{Unpriceable lines} (unmatched description, malformed date, unknown facility or tier) carry the
billed amount through as a \emph{provisional} expected total at reduced confidence --- never dropped, never
assigned an invented rate. Assumed better to state uncertainty than to publish a confident wrong figure.
\item \textbf{Confidence} is a deterministic audit-policy score reflecting evidence quality, not a
probability, and is reported as such.
\end{itemize}

\vspace{-3pt}
\section*{\normalsize 2.\; Clauses that read two ways, and the reading taken}
\vspace{-4pt}
\begin{itemize}
\item \textbf{Duplicate-billing scope (H4 11.3).} The clause bars billing a service twice for the same
patient and date, but cites a section listing only 18 capped services. \emph{Reading A:} it covers only
those 18. \emph{Reading B:} it is free-standing and covers every service. \textbf{Took B.} Consequence:
5 duplicate lines across 4 invoices, all on uncapped services, are flagged that Reading A would leave clean.
\item \textbf{Threshold-premium base.} Whether the premium applies to every unit on a day the threshold is
crossed, or only to the units above it. \textbf{Took ``all units'',} once the daily aggregate breaches the
threshold. The contract text cannot settle this either way; the narrower reading stays defensible.
\item \textbf{Daily-cap allocation.} Contracts state per-patient per-day maxima but not how to allocate them
across several line items. \textbf{Took ascending \texttt{line\_id} order}, zeroing the billable quantity of
lines beyond the cap.
\item \textbf{Volume-discount scope.} ``Aggregated across all patients'' does not say how to group.
\textbf{Took hospital-wide across the whole term}, counted in service-date then \texttt{line\_id} order, the
discount applying strictly \emph{after} the line that crosses the threshold; where two tiers are met, the
deeper applies.
\item \textbf{Premium vs.\ uplift ordering (H5 3.1(d)).} The clause groups ``any premium or uplift'' into one
stage without ordering them. \textbf{Took Section 5 premium, then Section 6 uplift}, following H1. The
alternative differs only by rounding, and only on lines where both fire.
\item \textbf{Which reused invoice ID is the erroneous one.} No contract says. \textbf{Took the later
occurrence}, recovering the transaction group from the line-ID prefix and mapping groups chronologically onto
invoice occurrences. Validated against H1's labels --- a calibration choice, disclosed as such rather than
presented as contract-derived.
\item \textbf{Two sources for the same rate tables (H5).} The tables ship inline in the agreement and again as
a separate PDF, with no stated precedence. Rather than pick one, both were extracted and compared
value-by-value before any pricing code was written: all 84 services agree exactly on unit basis, base rate,
daily cap and all six multipliers, so the Markdown is parsed as the single source. Had they disagreed, the
disagreement itself --- not a choice between them --- would have been the finding to report.
\end{itemize}

\vspace{-3pt}
\section*{\normalsize 3.\; Ambiguities not resolved, and what was done about each}
\vspace{-4pt}
\begin{itemize}
\item \textbf{Daily-cap corrected totals (H1, 4 invoices).} The literal cap reading does not reproduce the
labelled totals. In three of the four, the implied ``correct'' billable quantity is exactly 3 units
regardless of the service's stated cap (4, 12, 8) or the quantity billed (9, 14, 11); the fourth
(\texttt{INV-H1-000049}) carries a simultaneous premium error on the same line, so its implied quantity (9)
is confounded and not independent evidence. Cross-invoice cumulative counts and unit-basis conversion were
both ruled out, and Section 8 carries no interpretive clause. \emph{Done:} kept the contract-faithful
reading and disclosed the residual rather than hard-coding the labels. Most likely an artifact of how the
synthetic errors were injected, not a rule an auditor could recover.
\item \textbf{Descriptions naming no clinical specialty (H1, \texttt{INV-H1-000236}).} ``Fract Outpatient
Radiotherapy'' matches confidently to the sole Outpatient-prefixed candidate, where the label reads
\texttt{unknown\_service}. Two general fixes were tested and \emph{rejected on evidence}: raising the score
and gap thresholds did not change this match while breaking two correct ones; a missing-anchor penalty fixed
it but cost 47 false positives (precision $1.000\to0.552$). \emph{Done:} accepted as a disclosed limitation
--- one invoice in 913 is not worth 47 wrong accusations.
\item \textbf{Bundled rates as secondary price evidence.} The matcher's tie-break ignores bundle-substituted
rates in H1 and H4, so bundle-linked ambiguous descriptions are systematically under-resolved. \emph{Done:}
left them in the review queue under ``flag, don't guess''; closed for H5, whose plausible-price set models
both the multipliers and the bundled rate.
\item \textbf{No ground truth beyond Hospital 1.} H4 and H5 inherit H1's calibration rather than being
verified against labels, so their confidence figures are carried over, not measured. \emph{Done:} validated
H5 internally instead --- all 974 invoices it reports clean reconcile exactly against billed totals, and that
holds at 100\% independently within each of the three facilities and each of the three plan tiers, the check
that would fail first if a multiplier column were mismapped. This confirms the engine reproduces the
contract; it cannot detect a misreading of the contract itself.
\end{itemize}

\end{document}
